\section{后验坍缩与表示能力的失衡}
\label{sec:14-Deep-Learning/06-Variational-Autoencoders/03}

一个训练了两天的 VAE，各项指标都对：ELBO 单调上升后趋于平稳，没有 NaN，梯度范数正常，从先验采样出来的图片虽然糊，但确实像手写数字。你打算做一件很自然的事——拿一张真实图片编码成 \(z\)，然后把 \(z\) 的第三维从 \(-2\) 扫到 \(+2\)，看看这一维学到了什么。

输出一动不动。换第五维，还是不动。把整个 \(z\) 换成完全不同的另一个向量，解码结果几乎一样。再看编码端：不同的 \(x\) 输进去，\(\mu_\phi(x)\) 全都贴在零附近，\(\sigma_\phi(x)\) 全都贴在一附近。打印 KL 项，是 \(0.003\)。

这不是 bug。模型学到的东西是：\(q_\phi(z\given x)\approx p(z)\) 对所有 \(x\) 成立。近似后验完全不看输入，隐变量不携带任何关于 \(x\) 的信息，decoder 退化成一个无条件生成器。这就是后验坍缩。

上一节说过梯度接通不等于隐变量有用。现在把那句话展开：坍缩不是优化没跑好，它是我们写下的目标函数\emph{允许}甚至\emph{奖励}的一个解。

\subsection{把 ELBO 当作一份合同来读}

回到\hyperref[sec:14-Deep-Learning/06-Variational-Autoencoders/01]{``VAE 把不可积后验变成可训练下界''}里那个分解：

\[
\mathcal L(\theta,\phi;x) = \E_{q_\phi(z\given x)}\big[\log p_\theta(x\given z)\big] - \KL\big(q_\phi(z\given x)\,\Vert\, p(z)\big).
\]

优化器读到的是一份合同：第一项越大越好，第二项越小越好。它不知道你想要一个``有意义的隐表示''，它只知道这两个数。

现在设想 decoder 的能力足够强，强到不需要 \(z\) 也能把数据拟合得相当好——最典型的是自回归 decoder，它按像素（或按 token）逐个生成，每一步都能看到前面已经生成的全部内容。这类 decoder 本身就是一个完整的密度模型，把 \(z\) 整根拔掉，重构项也掉不了多少。

于是优化器面前有两个方案。方案甲：让 encoder 认真编码，\(z\) 携带 \(x\) 的信息，重构项拿到 \(+\delta\) 的提升，代价是 \(q_\phi(z\given x)\) 必须偏离先验，KL 项付出 \(R\)。方案乙：让 \(q_\phi\) 退回先验，KL 项精确付零，重构项少拿 \(\delta\)。只要 \(\delta < R\)，方案乙的 ELBO 更高。

强 decoder 让 \(\delta\) 变得很小，坍缩就成了\emph{全局最优}，不是局部陷阱。模型没有算错任何一步，是我们写的目标把这个解排在了前面。

这件事值得单独记住，因为它是全书反复出现的那条判断的一个干净实例：优化器不会理解你的意图，它会精确地钻目标函数的空子。你以为 KL 项在``正则化潜空间''，它同时也在给``不使用潜空间''发奖金。凡是把两个诉求写进同一个和式的目标，都要问一句：其中一项归零时，另一项能不能独自把总分撑住？

\subsection{KL 项是隐变量携带信息量的上界}

上面用 \(\delta\) 和 \(R\) 记了两个量，这个记号不是随手取的。把 ELBO 在数据分布上取平均，KL 项那一半可以精确拆开。记聚合后验为 \(q_\phi(z)=\E_{p(x)}[q_\phi(z\given x)]\)，则

\[
\E_{p(x)}\big[\KL(q_\phi(z\given x)\Vert p(z))\big]
= I_q(X;Z) + \KL\big(q_\phi(z)\,\Vert\, p(z)\big),
\]

其中 \(I_q(X;Z)\) 是在联合分布 \(p(x)q_\phi(z\given x)\) 下 \(X\) 与 \(Z\) 的互信息。第二项非负，所以

\[
I_q(X;Z) \le \E_{p(x)}\big[\KL(q_\phi(z\given x)\Vert p(z))\big].
\]

平均 KL 是隐变量所能携带信息量的上界。KL 归零，互信息必然归零——隐变量与输入统计独立，这跟``\(z\) 改变而输出不变''的现象是同一件事的两种说法。互信息本身的含义在\hyperref[sec:07-Information-Theory/02-Joint-Information/02]{``互信息是依赖，也是对不确定性的减少''}里有交代。

这个读法把 ELBO 变成了一个率失真问题：KL 项是``码率''，是你允许 \(z\) 这条信道传输的比特数；重构项是``失真''，是给定这些比特后还原 \(x\) 的能力。在固定的信道预算下，两者是零和的——想少花码率就得忍受失真，想降低失真就得付码率。ELBO 恰好是把两者以权重 \(1{:}1\) 相加。同一条曲线上的不同点对应不同的取舍，而 ELBO 只挑其中一个特定的点，那个点未必是你想要的表示。\hyperref[sec:13-Machine-Learning/11-Information-Theoretic-Learning/04]{``信息瓶颈在压缩与预测之间权衡''}讨论的是同一类结构：一个压缩项对着一个预测项，权重决定落在曲线的哪一段。

有了这个视角，``强 decoder 导致坍缩''就有了信息论的说法：decoder 越强，用零码率能达到的失真水平就越低，率失真曲线在零码率端越平，于是买第一个比特的性价比越差。

\subsection{诊断：看每一维的 KL，不看总和}

总 KL 是一个标量，它把所有维度的行为压成一个数，遮蔽了最需要看的结构。真正有用的诊断量是逐维的平均 KL：

\[
k_j = \E_{p(x)}\Big[\KL\big(q_\phi(z_j\given x)\,\Vert\, p(z_j)\big)\Big],
\qquad j=1,\dots,d .
\]

对角高斯下每一维都有闭式，算它几乎不花代价。把 \(k_1,\dots,k_d\) 画成柱状图，健康与坍缩的区别一眼可辨。

\begin{center}
\begin{tikzpicture}[>=stealth,x=1cm,y=1cm]
  \begin{scope}
    \draw[->] (0,0) -- (4.9,0) node[right,font=\scriptsize] {维度};
    \draw[->] (0,0) -- (0,3.1) node[above,font=\scriptsize] {每维平均 KL};
    \foreach \yv in {1,2}{\draw (0,\yv) -- (-0.12,\yv) node[left,font=\scriptsize] {\(\yv\)};}
    \foreach \i/\v in {1/2.45,2/2.10,3/1.95,4/1.70,5/1.55,6/1.30,7/1.15,8/0.90,9/0.70,10/0.55,11/0.35,12/0.22}
      \fill[blue!55!black,fill opacity=0.8] ({0.20+0.36*(\i-1)},0) rectangle ({0.20+0.36*(\i-1)+0.26},\v);
    \draw[dashed,black!50] (0,0.15) -- (4.7,0.15);
    \node[font=\small] at (2.35,-0.75) {健康模型};
  \end{scope}
  \begin{scope}[shift={(6.3,0)}]
    \draw[->] (0,0) -- (4.9,0) node[right,font=\scriptsize] {维度};
    \draw[->] (0,0) -- (0,3.1);
    \foreach \yv in {1,2}{\draw (0,\yv) -- (-0.12,\yv);}
    \foreach \i/\v in {1/2.55,2/1.05,3/0.09,4/0.06,5/0.05,6/0.05,7/0.04,8/0.04,9/0.03,10/0.03,11/0.02,12/0.02}
      \fill[red!60!black,fill opacity=0.8] ({0.20+0.36*(\i-1)},0) rectangle ({0.20+0.36*(\i-1)+0.26},\v);
    \draw[dashed,black!50] (0,0.15) -- (4.7,0.15);
    \node[font=\scriptsize,black!60,above] at (3.3,0.22) {活跃阈值};
    \node[font=\small] at (2.35,-0.75) {坍缩模型};
  \end{scope}
\end{tikzpicture}

\emph{总 KL 相近的两个模型，逐维分布可以完全不同：健康模型把码率摊在许多维上，坍缩模型只剩一两维在贴地线之上。诊断要看这张图的形状，不要看它的面积。}
\end{center}

图里那条虚线是活跃阈值，通常取 \(0.01\) 到 \(0.1\) 之间的一个小数，超过它的维度记为活跃单元。活跃单元数随训练进程的曲线比 ELBO 曲线信息量大得多：它掉到个位数并且不再回升，基本可以判定坍缩。另一个等价的做法是看 \(\mu_\phi(x)\) 在数据集上的逐维方差——某一维在所有输入上都输出同一个数，这一维就没在工作。这些阈值都是经验规律，不同数据集和不同潜维度下要重新校准，没有理论给出的分界点。

顺带一提，右图那种``只有一两维活跃''的形状本身不一定是坏事。如果数据的内在维度确实只有两维，模型自动关掉多余的维度反而是对的。判断坍缩要结合重构质量：活跃维度少\emph{且}重建明显差于 decoder 的能力上限，才是问题。

\subsection{四种缓解手段，动的不是同一个东西}

知道坍缩来自目标函数的结构之后，各种缓解手段的差别就清楚了——它们分别在改优化路径、改目标、改约束、改模型。

KL 退火在训练早期给 KL 项一个从 \(0\) 升到 \(1\) 的权重 \(\beta_t\)。前期 KL 几乎免费，encoder 被鼓励大胆编码；等 \(z\) 已经承担了信息，再把 KL 的价钱慢慢提上来。它改变的是\emph{优化路径}，不是最优解——权重回到 \(1\) 之后，目标还是原来那个 ELBO，理论上的最优解一点没变。它能起作用，依赖的是``优化器已经走到一个使用 \(z\) 的盆地里，不会再翻回去''，这是一条经验规律，训练久了退回坍缩的情况确实发生过。

\(\beta\)-VAE 把 KL 项永久加权成 \(\beta\cdot\KL\)。取 \(\beta<1\) 直接削弱对隐变量的惩罚，坍缩自然缓解。代价必须说明白：加权之后的量\emph{不再是 \(\log p_\theta(x)\) 的下界}——Jensen 不等式给出的是权重为 \(1\) 的那个式子，改了权重，不等式就失效了。你仍然在优化一个正当的率失真目标，只是不能再报告``我们最大化了似然下界''，也不能拿这个数跟别人的 ELBO 比。取 \(\beta>1\) 是另一个方向的用法，为了换取更强的解耦，代价是重构变差。

Free bits 给每一维的 KL 设一个地板：把上面那个逐维量 \(k_j\) 换成 \(\max(\lambda, k_j)\) 再求和，\(\lambda\) 是每维允许免费使用的信息量。低于地板的维度不受惩罚，优化器没有动力把它压到零；高于地板的部分照常付费。它改的是\emph{约束的形状}，比全局退火更有针对性，因为惩罚是逐维施加的。代价是 \(\max\) 引入了不可微点，且 \(\lambda\) 需要调。

削弱 decoder 则从模型容量入手：不用自回归 decoder，或者给 decoder 加 dropout、限制感受野，让它没有 \(z\) 就真的做不好。这是唯一一个不改目标函数的办法——它改的是那个 \(\delta\)，让``使用隐变量''重新变得划算。代价是模型上限降低了。

摆在一起看，前三种是在跟目标函数讨价还价，第四种是在改变博弈的收益矩阵。哪种合适取决于你要的是什么：如果目的是拿到可解释的隐表示，改目标是可以接受的；如果目的是密度估计并且要报告似然，那么改目标就等于改了你在报告的东西。

\subsection{KL 不为零，也不等于隐变量有用}

最后一条提醒是诊断层面的，容易被跳过。上面所有手段都能把 KL 从零拉起来，但 KL 大于零只说明 \(z\) 与 \(x\) 之间有\emph{某种}统计依赖，没说这依赖是什么。

一个常见的失望情形：free bits 把每维 KL 稳稳顶在 \(0.5\) 以上，活跃单元数看着很健康，可你去扫描每一维，发现它们编码的是图像整体亮度、背景灰度、笔画粗细这类低层统计量——这些量最容易编码，付出的 KL 最少而换来的重构收益最直接。真正想要的语义因素（数字的类别、书写风格）仍然由 decoder 独自建模。模型精确地满足了你写下的约束，而不是你心里想的那件事。

所以逐维 KL 直方图是必要条件不是充分条件。要判断表示是否有用，得做别的检查：把 \(z\) 拿去做下游分类，看线性探针的准确率；固定 \(z\) 的一维做遍历，看输出的变化是否语义连贯；或者比较 \(q_\phi(z\given x)\) 与打乱标签后的表现。这些检查跟 ELBO 高低只有松散的相关，必须单独做。

坍缩这个现象本身指向一个更大的问题：VAE 只是把``最小化某个散度''的思路实例化成了 ELBO，而散度的选择、下界的松紧、目标允许的退化解，都是这一次实例化带来的后果，不是生成建模的必然。下一个单元把 VAE 与 GAN、流模型、扩散模型放在同一个框架下对照，问的正是这个问题：每一类模型实际在最小化哪个散度，那个选择又各自决定了它们会在什么地方失败。
